\section{SNCASE Mistral}

The SNCASE (Societé nationale des constructions aéronautiques du Sud-Est) SE.535 Mistral is a day fighter and fighter-bomber. It is named for the strong wind the occurs frequently in Provence.

The Mistral is a development of the SNCASE Vampire FB.5, but fitted with the more powerful Nene 104B motor (license-built by Hispano-Suiza) providing 5,000 lb of thrust compared to 3,000 lb for the Goblin 2 engine of the FB.5 and 3,500 lb for the Goblin 3 engine of the FB.6. The Nene engine had been trialed in the Vampire F.2 with its need for greater airflow satisfied by additional “elephant ears” inlets on the upper side of the fuselage behind the cockpit. However, these intakes caused handling problems close to the critical Mach number. In the DHA Vampire F.30, this was mitigated by moving the addition intakes to the underside of the fuselage. In contrast, the Mistral omitted them completely and instead opted for redesigned and enlarged wing-root inlets, and as a consequence had much more benign characteristics at high speed. The Mistral also featured an ejection seat and air conditioning to facilitate its use in the French colonies in Africa.

It is not clear how much de Havilland was involved in the development of the Mistral, but nevertheless
early prototypes of the Mistral are sometimes referred to as the Vampire FB.53.

The Mistral served in the French AA from 1953 until at least 1961 and saw combat in the colonial wars in Tunisia and Algeria.

Its gun armament is four 20 mm Hispano cannons with 150 rounds, as in the Vampire. Air-to-ground ordnance included rockets, HE bombs, and napalm bombs.

An ADC is provided for the:
\begin{adclist}
    \adcitem{Mistral}
\end{adclist}

See also:
\begin{typelist}
    \typeitem{de Havilland Vampire}
    \typeitem{DHA Vampire}
    \typeitem{SNCASE Vampire}
\end{typelist}